% !TEX program = xelatex
% 공통수학2 · Ⅰ. 도형의 방정식 · 1. 평면좌표와 직선의 방정식 · 중단원 마무리 (8차시)
% 학생용 활동지 (답안 없음)

\documentclass[11pt,a4paper]{article}

\usepackage{kotex}
\usepackage{iftex}
\ifPDFTeX\else
  \setmainhangulfont{Noto Serif CJK KR}
  \setsanshangulfont{Noto Sans CJK KR}
\fi

\usepackage[margin=2.1cm]{geometry}
\usepackage{parskip}
\usepackage{enumitem}
\usepackage{array}
\usepackage{tabularx}
\usepackage{xcolor}
\usepackage{amssymb}
\usepackage{tikz}
\usepackage[most]{tcolorbox}
\usepackage{fancyhdr}
\usepackage{titlesec}

\definecolor{goldc}{HTML}{B07C22}
\definecolor{goldstrong}{HTML}{8A5F16}
\definecolor{indigoc}{HTML}{2B3B57}
\definecolor{indigosoft}{HTML}{6E82A6}
\definecolor{linec}{HTML}{D6DACB}
\definecolor{surface2}{HTML}{E4E8DC}
\definecolor{writeline}{HTML}{C7CCBA}
\definecolor{inksoft}{HTML}{52604F}

\titleformat{\section}{\normalfont\sffamily\bfseries\large\color{indigoc}}{\thesection}{0.6em}{}
\titlespacing{\section}{0pt}{14pt}{6pt}

\newtcolorbox{goalbox}{enhanced, breakable, colback=white, colframe=goldc, boxrule=0.5pt, leftrule=3pt, arc=2pt, left=10pt, right=10pt, top=8pt, bottom=8pt}
\newtcolorbox{sheetbox}{enhanced, breakable, colback=white, colframe=linec, boxrule=0.6pt, arc=3pt, left=11pt, right=11pt, top=9pt, bottom=9pt}

\newcommand{\writespace}[1]{%
  \par\nobreak\vspace{3pt}
  \foreach \n in {1,...,#1} {%
    \noindent\textcolor{writeline}{\rule{\linewidth}{0.4pt}}\par\vspace{0.62cm}%
  }
}
\newcommand{\fillblank}[1]{\makebox[#1]{\hrulefill}}

\pagestyle{fancy}
\fancyhf{}
\renewcommand{\headrulewidth}{0pt}
\fancyfoot[L]{\footnotesize\color{inksoft} 공통수학2 · 2026학년도 2학기 · 지평선고등학교}
\fancyfoot[R]{\footnotesize\color{inksoft} 다음 시간 --- 2. 원의 방정식}

\begin{document}

\noindent
\begin{minipage}[b]{0.62\linewidth}
  {\sffamily\footnotesize\color{indigosoft} 공통수학2 · Ⅰ. 도형의 방정식}\\[2pt]
  {\sffamily\bfseries\Large 1. 평면좌표와 직선의 방정식 \textperiodcentered\ 중단원 마무리}
\end{minipage}\hfill
\begin{minipage}[b]{0.36\linewidth}
  \raggedleft\small\color{inksoft}
  반\ \fillblank{1.6cm} \quad 번호\ \fillblank{1.6cm}\\[4pt]
  이름\ \fillblank{3.6cm}
\end{minipage}

\vspace{4pt}
\noindent{\color{goldc}\rule{\linewidth}{2.2pt}}
\vspace{10pt}

\begin{goalbox}
\textbf{오늘의 목표}\\
1$\sim$7차시에서 배운 개념(내분, 평행·수직 조건, 점과 직선 사이의 거리)을 스스로 정리하고 점검한다.
\end{goalbox}

\section{개념 지도 채우기}

\begin{sheetbox}
\begin{tabularx}{\linewidth}{@{}>{\bfseries\small}p{4.4cm}X@{}}
\toprule
\textbf{개념} & \textbf{공식} \\
\midrule
수직선 위 두 점 사이의 거리 & \fillblank{6.5cm} \\[6pt]
수직선 위 내분점 & \fillblank{6.5cm} \\[6pt]
좌표평면 위 두 점 사이의 거리 & \fillblank{6.5cm} \\[6pt]
좌표평면 위 내분점 & \fillblank{6.5cm} \\[6pt]
삼각형의 무게중심 & \fillblank{6.5cm} \\[6pt]
두 직선의 평행 조건 & \fillblank{6.5cm} \\[6pt]
두 직선의 수직 조건 & \fillblank{6.5cm} \\[6pt]
점과 직선 사이의 거리 & \fillblank{6.5cm} \\
\bottomrule
\end{tabularx}
\end{sheetbox}

\section{기본 문제}

\begin{sheetbox}
\textbf{1.} 두 점 A$(3)$, B$(-7)$ 사이의 거리와 두 점 A$(-5,1)$, B$(-2,6)$ 사이의 거리를 각각 구하시오.
\writespace{2}

\textbf{2.} 두 점 A$(6,2)$와 B$(-3,-1)$에 대하여 선분 AB를 $1:2$로 내분하는 점의 좌표를 구하시오.
\writespace{2}

\textbf{3.} 점 $(1,-3)$을 지나고 직선 $y=2x-1$에 수직인 직선의 방정식을 구하시오.
\writespace{2}

\textbf{4.} 점 $(4,2)$와 직선 $3x-4y+11=0$ 사이의 거리를 구하시오.
\writespace{2}
\end{sheetbox}

\section{표준 \& 도전 문제}

\begin{sheetbox}
\textbf{5.} 세 점 A$(-1,2)$, B$(3,5)$, C$(a,b)$를 꼭짓점으로 하는 삼각형 ABC의 무게중심의 좌표가 $(1,1)$일 때, $a+b$의 값을 구하시오.
\writespace{3}

\textbf{6.} 평행한 두 직선 $2x+3y-1=0$과 $2x+3y-k=0$ 사이의 거리가 $\sqrt{13}$일 때, 상수 $k$의 값을 모두 구하시오.
\writespace{3}

\textbf{7. (서술형)} 두 직선 $2x-3y+4=0$과 $3x+y-5=0$의 교점을 지나고 직선 $6x+3y+1=0$과 평행한 직선의 방정식을 구하는 풀이 과정과 답을 쓰시오.
\writespace{4}
\end{sheetbox}

\section{마무리 성찰 --- 나의 성취 수준 점검}

\begin{sheetbox}
아래 항목 중 자신 있게 설명할 수 있는 것에 $\checkmark$ 표시해 보자.

\begin{itemize}[leftmargin=1.6em, itemsep=3pt]
  \item[$\square$] 수직선·좌표평면 위의 내분점과 중점을 구할 수 있다.
  \item[$\square$] 삼각형의 무게중심을 구할 수 있다.
  \item[$\square$] 두 직선의 평행·수직 조건을 설명할 수 있다.
  \item[$\square$] 점과 직선 사이의 거리를 구할 수 있다.
  \item[$\square$] 평행선 사이의 거리와 삼각형의 넓이를 구할 수 있다.
\end{itemize}

\vspace{6pt}
{\footnotesize\color{inksoft} 가장 보충이 필요한 부분과 그 이유를 적어 보자.}
\writespace{2}
\end{sheetbox}

\end{document}
